% Problem record op_787284adc73c8f63. This TeX file is derived from
% database/problems_json/op_787284adc73c8f63.json by scripts/sync-tex.mjs; edit the JSON record
% and rerun the script rather than editing this file.
\section{Exact Gilbert--Varshamov distance for rate-one-half quantum LDPC codes}
\label{sec:op_787284adc73c8f63}

\paragraph{Problem.}
Can qubit stabilizer codes with uniformly sparse checks attain the quantum Gilbert--Varshamov relative distance at rate $1/2$? Let $\delta_*$ be the unique root in $(0,1/2)$ of $h_2(\delta_*)+\delta_*\log_2 3=1/2$, where $h_2$ is binary entropy. For a code $[[n,k_n,d_n]]$, the distance $d_n$ is the smallest weight of a Pauli that preserves the codespace and acts nontrivially on it.

Seek a family with $n\to\infty$ and
\begin{equation}
\liminf_n k_n/n\geq 1/2,\qquad
\liminf_n d_n/n\geq \delta_*.
\label{eq:gvq-rate-distance}
\end{equation}
There must be one constant $w$ for the whole family such that each code has stabilizer generators acting on at most $w$ qubits, with each qubit in at most $w$ of those generators. Does a family satisfying Eq.~\eqref{eq:gvq-rate-distance} exist?

\subsection*{Status}

Unsolved

\subsection*{Source}

Mittal, Srivastava, Tulsiani and Wootters ask in Section~1 whether quantum LDPC codes can reach the quantum Gilbert--Varshamov bound. Equation~\eqref{eq:gvq-rate-distance} specializes their question to binary codes and rate $1/2$ with one fixed bound on generator weight and incidence \sourcecite{ref:gvq-mstw}{MSTW26}.

\subsection*{Progress}

\begin{itemize}
  \item Mittal, Srivastava, Tulsiani and Wootters show that for each fixed $\varepsilon>0$, sparse qubit stabilizer codes of rate $1/2$ can have normalized distance at least $\delta_*-\varepsilon$. The required check-weight bound depends on $\varepsilon$, so choosing successively smaller $\varepsilon$ gives no single bound for the endpoint \sourcecite{ref:gvq-mstw}{MSTW26}.

  \item Kasai's Corollary~5.2 attains the CSS Gilbert--Varshamov benchmark for a collection of finite-degree ensembles. Section~2.4 explains that sparsity of the extended system does not establish bounded-weight measurements of the compressed physical syndromes \sourcecite{ref:gvq-kasai}{Kas26}.
\end{itemize}

\subsection*{References}

\begin{enumerate}[label={},leftmargin=4.8em,labelsep=0.5em]
  \footnotesize
  \item[\textup{[MSTW26]}]\label{ref:gvq-mstw}
  T. Mittal, S. Srivastava, M. Tulsiani, and M. Wootters, ``Random Quantum LDPC Codes Approaching the Gilbert--Varshamov Bound,'' arXiv preprint (2026). \href{https://arxiv.org/abs/2610.02648v1}{arXiv:2610.02648v1}.

  \item[\textup{[Kas26]}]\label{ref:gvq-kasai}
  K. Kasai, ``Finite-Degree Quantum LDPC Codes Reaching the Gilbert--Varshamov Bound,'' arXiv preprint (2026). \href{https://arxiv.org/abs/2603.24588v3}{arXiv:2603.24588v3}.
\end{enumerate}

\subsection*{Comment}

The locality constant must stay fixed as $n$ grows, including at exact asymptotic equality. The related \href{https://qiqc-op.com/problem/op_4830398d0c2feb5f/}{Pauli hashing-bound question} uses stochastic errors. The \href{https://qiqc-op.com/problem/op_edcb345719181830/}{qubit distance-bound question} asks for tightness of an upper bound across rates. The catalog and literature were checked through November 26, 2026.

\subsection*{Field}

Quantum Error Correction

\subsection*{Topic}

Quantum LDPC codes; Quantum coding theory

\subsection*{ID}

\texttt{op\_787284adc73c8f63}
